\documentclass[aspectratio=169,10pt]{beamer}
\usepackage[utf8]{inputenc}
\usepackage[T1]{fontenc}
\usepackage{lmodern}
\usepackage{amsmath,amssymb}
\usepackage{tikz}
\usetikzlibrary{arrows.meta,positioning,fit,backgrounds,calc,shadows,shapes.geometric}
\usepackage{fontawesome5}
\usepackage{tcolorbox}
\tcbuselibrary{skins,raster}

% Theme Colors
\definecolor{lavenderDark}{RGB}{76, 51, 122}
\definecolor{lavenderMedium}{RGB}{124, 98, 173}
\definecolor{lavenderLight}{RGB}{243, 239, 252}
\definecolor{lavenderAccent}{RGB}{150, 120, 205}
\definecolor{prismDark}{RGB}{23, 32, 42}
\definecolor{prismSlate}{RGB}{85, 90, 105}
\definecolor{prismLight}{RGB}{248, 247, 252}
\definecolor{prismTeal}{RGB}{22, 138, 114}
\definecolor{prismBlue}{RGB}{124, 98, 173}

\setbeamercolor{frametitle}{bg=lavenderDark,fg=white}
\setbeamercolor{author in head/foot}{bg=lavenderDark, fg=white}
\setbeamercolor{title in head/foot}{bg=lavenderMedium, fg=white}
\setbeamercolor{date in head/foot}{bg=lavenderDark, fg=white}
\setbeamertemplate{navigation symbols}{}

% Hierarchical Algorithm List Environments matching Screenshot 2 structure
\newenvironment{algolist}{%
  \begin{list}{}{%
    \setlength{\labelwidth}{1.9em}%
    \setlength{\leftmargin}{2.1em}%
    \setlength{\itemindent}{0pt}%
    \setlength{\labelsep}{0.45em}%
    \setlength{\itemsep}{0.28em}%
    \setlength{\parsep}{0pt}%
    \setlength{\topsep}{0.1em}%
    \renewcommand{\makelabel}[1]{\textbf{\color{lavenderDark}##1}}%
  }\small%
}{\end{list}}

\newenvironment{algosublist}{%
  \begin{list}{}{%
    \setlength{\labelwidth}{2.4em}%
    \setlength{\leftmargin}{2.8em}%
    \setlength{\itemindent}{0pt}%
    \setlength{\labelsep}{0.4em}%
    \setlength{\itemsep}{0.14em}%
    \setlength{\parsep}{0pt}%
    \setlength{\topsep}{0.1em}%
    \renewcommand{\makelabel}[1]{\textbf{\color{lavenderMedium}##1}}%
  }\small%
}{\end{list}}

\begin{document}

% -----------------------------------------------------------------------------
% STEP 1: PROMPT PARSING & TASK REPRESENTATION
% -----------------------------------------------------------------------------
\begin{frame}[t]{Algorithm: Step 1 -- Understand the User Prompt}
\vspace{-0.15em}
{\large\bfseries\color{lavenderDark}Step 1: Prompt Parsing and Structured Representation}
\vspace{0.15em}

\begin{algolist}
  \item[1.1] Ingest raw natural language prompt $\mathcal{T}_{\text{user}}$ via client interface and extract core objective.
  \item[1.2] Retrieve persistent session context and workspace history from browser IndexedDB (\texttt{PrismDB}):
    \begin{algosublist}
      \item[1.2.1] Load prior interaction turns, agent memory checkpoints, and active user configurations.
      \item[1.2.2] Resolve local file system references, target project repositories, and access tokens.
    \end{algosublist}
  \item[1.3] Construct strongly-typed operational execution tuple:
    \vspace{-0.2em}
    $$T = \langle \text{Goal}, \, \text{Constraints}, \, \text{Context}, \, \mathcal{M}_{\text{tools}}, \, \Phi_{\text{out}} \rangle$$
    \vspace{-0.2em}
    {\footnotesize\color{prismDark}%
      \textbf{\color{lavenderDark}Where:}\enspace \text{Goal}: Target objective $\;\vert\;$ \text{Constraints}: Safety \& SLA bounds $\;\vert\;$ \text{Context}: Session state in \texttt{PrismDB} $\;\vert\;$ $\mathcal{M}_{\text{tools}}$: Authorized MCP tools $\;\vert\;$ $\Phi_{\text{out}}$: Deliverable schema contract.}
    \vspace{0.15em}
    \begin{algosublist}
      \item[1.3.1] Parse explicit operational constraints (deadlines, schema specifications, output formats).
      \item[1.3.2] Bind authorized Model Context Protocol (MCP) tool capabilities and environment permissions.
    \end{algosublist}
  \item[1.4] Perform input boundary and security screening:
    \begin{algosublist}
      \item[1.4.1] Evaluate prompt against token limits and scan for prompt injection patterns.
      \item[1.4.2] Flag state-mutating actions requiring mandatory human-in-the-loop authorization.
    \end{algosublist}
\end{algolist}
\end{frame}

% -----------------------------------------------------------------------------
% STEP 2: INTENT CLASSIFICATION & CAPABILITY PROFILING
% -----------------------------------------------------------------------------
\begin{frame}[t]{Algorithm: Step 2 -- Classify the Request}
\vspace{-0.15em}
{\large\bfseries\color{lavenderDark}Step 2: Intent Classification and Capability Check}
\vspace{0.15em}

\begin{algolist}
  \item[2.1] Categorize user prompt into functional operational domains:
    \vspace{-0.2em}
    $$\text{Domain} \in \{\text{Productivity/Automation}, \, \text{Document Synthesis}, \, \text{Web Research}, \, \text{Code Engineering}\}$$
    \vspace{-0.2em}
    {\footnotesize\color{prismDark}%
      \textbf{\color{lavenderDark}Where:}\enspace \text{Productivity}: Communication \& planning $\;\vert\;$ \text{Synthesis}: Summarization $\;\vert\;$ \text{Research}: Web search \& scraping $\;\vert\;$ \text{Coding}: CLI \& code synthesis.}
    \vspace{0.15em}
    \begin{algosublist}
      \item[2.1.1] Compute dense semantic embeddings and syntactic keyword vectors for input prompt.
      \item[2.1.2] Route through multi-class classifier head $\hat{\mathbf{y}}_{\text{intent}}$ to determine target capability domain.
    \end{algosublist}
  \item[2.2] Profile required system capabilities and operational risk levels:
    \begin{algosublist}
      \item[2.2.1] Determine required subsystem resources (live web search, local filesystem I/O, terminal execution).
      \item[2.2.2] Enforce safety categorization: Read-Only (autonomous execution) vs. State-Mutating (approval-gated).
    \end{algosublist}
  \item[2.3] Select optimal execution pathway:
    \begin{algosublist}
      \item[2.3.1] Route single-step queries to direct, low-latency single-model inference pipeline.
      \item[2.3.2] Dispatch multi-step complex requests to DAG Planner agent for task decomposition.
    \end{algosublist}
\end{algolist}
\end{frame}

% -----------------------------------------------------------------------------
% STEP 3: DAG DECOMPOSITION & MODEL ROUTING
% -----------------------------------------------------------------------------
\begin{frame}[t]{Algorithm: Step 3 -- Plan and Route the Task}
\vspace{-0.15em}
{\large\bfseries\color{lavenderDark}Step 3: DAG Decomposition and Model Routing}
\vspace{0.15em}

\begin{algolist}
  \item[3.1] Deconstruct objective into Directed Acyclic Graph $\mathcal{G} = (V, E)$:
    \begin{algosublist}
      \item[3.1.1] Instantiate discrete task vertices $v_i \in V$ with atomic execution contracts and output targets.
      \item[3.1.2] \textbf{Data Dependencies as Directed Edges} ($e = \langle v_i, v_j \rangle \in E \iff \mathcal{O}_i \cap \mathcal{I}_j \neq \emptyset$): Task $v_j$ requires input artifacts produced by $v_i$ (e.g., synthesis requires scraped research).
      \item[3.1.3] \textbf{Execution Barrier:} Consumer $v_j$ remains blocked until producer $v_i$ delivers $\mathcal{O}_i$, guaranteeing tasks never run on missing or partial data.
    \end{algosublist}
  \item[3.2] Compute topological task execution schedule $\pi = \mathrm{TopologicalSort}(\mathcal{G})$:
    \begin{algosublist}
      \item[3.2.1] Identify independent vertices ($\text{in-degree} = 0$) for concurrent parallel dispatch ($K \le 3$).
      \item[3.2.2] Schedule downstream vertices dynamically upon upstream output event emission.
    \end{algosublist}
  \item[3.3] Solve multi-objective provider routing optimization for each sub-task $v_i$:
    \vspace{-0.2em}
    $$\text{Provider}^* = \arg\min_m \Big(\lambda_{\text{cost}} \, \text{Cost}_m + \lambda_{\text{latency}} \, \text{Latency}_m\Big) \quad \text{s.t. } \text{Quality}_m \ge \tau$$
    \vspace{-0.2em}
    {\footnotesize\color{prismDark}%
      \textbf{\color{lavenderDark}Where:}\enspace $\text{Provider}^*$: Optimal LLM backend $\;\vert\;$ $m$: Candidate model $\;\vert\;$ $\text{Cost}_m, \text{Latency}_m$: Normalized price \& delay $\;\vert\;$ $\lambda_{\text{cost}}, \lambda_{\text{latency}}$: Balancing weights $\;\vert\;$ $\tau$: Quality barrier.}
    \vspace{0.15em}
    \begin{algosublist}
      \item[3.3.1] Route high-throughput, latency-critical nodes to LPU backends (Groq / Llama-3).
      \item[3.3.2] Route deep reasoning and code generation tasks to frontier backends (OpenAI GPT-4o, Claude 3.5).
    \end{algosublist}
\end{algolist}
\end{frame}

% -----------------------------------------------------------------------------
% STEP 4: TASK ASSIGNMENT & AGENT SELECTION
% -----------------------------------------------------------------------------
\begin{frame}[t]{Algorithm: Step 4 -- Task Assignment and Agent Selection}
\vspace{-0.15em}
{\large\bfseries\color{lavenderDark}Step 4: Role-Based Task Assignment and Dispatching}
\vspace{0.15em}

\begin{algolist}
  \item[4.1] Query specialist agent registry and capability profiles:
    \vspace{-0.2em}
    $$\mathcal{A} = \{\text{Planner}, \, \text{Assistant}, \, \text{Researcher}, \, \text{Coder}, \, \text{Validator}\}$$
    \vspace{-0.2em}
    \begin{algosublist}
      \item[4.1.1] Evaluate persona system prompts, specialized instruction boundaries, and assigned MCP tools.
      \item[4.1.2] Verify agent worker availability and monitor concurrent thread utilization.
    \end{algosublist}
  \item[4.2] Execute capability matching algorithm for each sub-task node $v_i$:
    \vspace{-0.2em}
    $$\text{Agent}^* = \arg\max_{a \in \mathcal{A}} \; \mathrm{CosineSimilarity}\big(\mathbf{Req}(v_i), \, \mathbf{Cap}(a)\big)$$
    \vspace{-0.2em}
    {\footnotesize\color{prismDark}%
      \textbf{\color{lavenderDark}Where:}\enspace $\text{Agent}^*$: Assigned specialist persona $\;\vert\;$ $\mathcal{A}$: Swarm agent registry $\;\vert\;$ $\mathbf{Req}(v_i)$: Task requirement vector $\;\vert\;$ $\mathbf{Cap}(a)$: Agent capability embedding $\;\vert\;$ $\mathrm{CosineSimilarity}$: Semantic alignment.}
    \vspace{0.15em}
    \begin{algosublist}
      \item[4.2.1] Match information gathering requirements ($\text{Tag} \ni \texttt{web\_search}$) to Researcher agent.
      \item[4.2.2] Match software synthesis and execution ($\text{Tag} \ni \texttt{python, cli}$) to Coder agent.
      \item[4.2.3] Match quality assurance and policy compliance ($\text{Tag} \ni \texttt{audit, verify}$) to Validator agent.
    \end{algosublist}
  \item[4.3] Manage task lifecycle states in dispatcher queue:
    \begin{algosublist}
      \item[4.3.1] Place sub-tasks in \texttt{PENDING} state until all parent dependency inputs resolve.
      \item[4.3.2] Transition verified nodes to \texttt{READY} and dispatch execution payloads to assigned agents.
    \end{algosublist}
\end{algolist}
\end{frame}

% -----------------------------------------------------------------------------
% STEP 5: AGENT COMMUNICATION & DECISION MAKING
% -----------------------------------------------------------------------------
\begin{frame}[t]{Algorithm: Step 5 -- Agent Communication and Decision Cycle}
\vspace{-0.15em}
{\large\bfseries\color{lavenderDark}Step 5: Agent Decision Cycle and Inter-Agent Communication}
\vspace{0.15em}

\begin{algolist}
  \item[5.1] Execute autonomous Reasoning-Action (ReAct) cycle for each dispatched agent:
    \vspace{-0.2em}
    $$\text{State } s_t \;\longrightarrow\; \text{\textbf{Observe}} \;\longrightarrow\; \text{\textbf{Reason}} \;\longrightarrow\; \text{\textbf{Act}} \;\longrightarrow\; \text{\textbf{Reflect}}$$
    \vspace{-0.2em}
    {\footnotesize\color{prismDark}%
      \textbf{\color{lavenderDark}Where:}\enspace $s_t$: Current execution state at turn $t$ $\;\vert\;$ \textbf{Observe}: Ingest context \& parent outputs $\mathcal{O}_i$ $\;\vert\;$ \textbf{Reason}: Plan via CoT $\;\vert\;$ \textbf{Act}: Invoke MCP tools $\;\vert\;$ \textbf{Reflect}: Self-correct on error.}
    \vspace{0.15em}
    \begin{algosublist}
      \item[5.1.1] \textbf{Observe:} Ingest assigned sub-task description, operational constraints, and parent inputs.
      \item[5.1.2] \textbf{Reason:} Formulate intermediate chain-of-thought rationale to plan next tool invocation.
      \item[5.1.3] \textbf{Act:} Dispatch structured JSON-RPC tool calls to external MCP servers or formulate text.
      \item[5.1.4] \textbf{Reflect:} Inspect execution exit codes and output schemas; initiate local correction on errors.
    \end{algosublist}
  \item[5.2] Transmit intermediate artifacts via DAG edge message passing:
    \begin{algosublist}
      \item[5.2.1] Route validated producer artifact $\mathcal{O}_i$ directly into consumer input buffer $\mathcal{I}_j$.
      \item[5.2.2] Enforce strict artifact isolation to prevent cross-task context contamination and token inflation.
    \end{algosublist}
  \item[5.3] Synchronize shared state via local blackboard memory:
    \begin{algosublist}
      \item[5.3.1] Commit global project state, conversation turns, and file handles to IndexedDB (\texttt{PrismDB}).
      \item[5.3.2] Dispatch event-driven signals upon node resolution to awaken dependent downstream workers.
    \end{algosublist}
\end{algolist}
\end{frame}

% -----------------------------------------------------------------------------
% STEP 6: SPECIALIST WORKFLOWS: PLANNER & ASSISTANT
% -----------------------------------------------------------------------------
\begin{frame}[t]{Algorithm: Step 6 -- Specialist Workflows: Planner \& Assistant}
\vspace{-0.15em}
{\large\bfseries\color{lavenderDark}Step 6: Specialist Execution Workflows -- Planner \& Assistant}
\vspace{0.15em}

\begin{algolist}
  \item[6.1] Execute Planner Agent workflow (System Architect \& Coordinator):
    \begin{algosublist}
      \item[6.1.1] Ingest user objective string and conversation context from \texttt{PrismDB}.
      \item[6.1.2] Decompose objective into atomic independent and dependent sub-goals.
      \item[6.1.3] Construct DAG graph vertices $V$, dependency edges $E$, and capability tags.
      \item[6.1.4] Solve provider routing optimization and enqueue topological schedule into Dispatcher.
    \end{algosublist}
  \item[6.2] Execute Assistant Agent workflow (Productivity \& Communication):
    \begin{algosublist}
      \item[6.2.1] Retrieve user communication preferences, profiles, and reference documents from storage.
      \item[6.2.2] Formulate structured response drafts, templates, or executive summaries.
      \item[6.2.3] Interface with Google Workspace MCP tools (Gmail draft creation, Drive document retrieval).
      \item[6.2.4] Emit verified communication draft; enforce approval gate prior to external message transmission.
    \end{algosublist}
\end{algolist}
\end{frame}

% -----------------------------------------------------------------------------
% STEP 7: SPECIALIST WORKFLOWS: RESEARCHER, CODER & VALIDATOR
% -----------------------------------------------------------------------------
\begin{frame}[t]{Algorithm: Step 7 -- Specialist Workflows: Researcher, Coder \& Validator}
\vspace{-0.15em}
{\large\bfseries\color{lavenderDark}Step 7: Specialist Execution Workflows -- Researcher, Coder \& Validator}
\vspace{0.15em}

\begin{algolist}
  \item[7.1] Execute Researcher Agent workflow (Information Retrieval \& Synthesis):
    \begin{algosublist}
      \item[7.1.1] Formulate optimized web search queries and invoke Browser MCP server endpoints.
      \item[7.1.2] Scrape target web documents, parse DOM trees, and extract relevant contextual passages.
      \item[7.1.3] Deduplicate findings and compile structured Markdown report with verified source citations.
    \end{algosublist}
  \item[7.2] Execute Coder Agent workflow (Software Synthesis \& Execution):
    \begin{algosublist}
      \item[7.2.1] Inspect local repository tree via Filesystem MCP and synthesize target source code files.
      \item[7.2.2] Execute scripts inside Sandboxed Terminal CLI and inspect exit codes and error streams.
      \item[7.2.3] Analyze stack traces and compiler lints to autonomously self-correct syntax and runtime bugs.
    \end{algosublist}
  \item[7.3] Execute Validator Agent workflow (Quality Auditing \& Safety Gate):
    \begin{algosublist}
      \item[7.3.1] Scan candidate outputs against safety policies, token invariants, and functional constraints.
      \item[7.3.2] Emit binary verification signal: \textbf{PASS} to release artifact, or \textbf{FAIL} with error diagnostics.
    \end{algosublist}
\end{algolist}
\end{frame}

% -----------------------------------------------------------------------------
% STEP 8: EXECUTE WITH MCP TOOLS
% -----------------------------------------------------------------------------
\begin{frame}[t]{Algorithm: Step 8 -- Standardized Tool Execution via MCP}
\vspace{-0.15em}
{\large\bfseries\color{lavenderDark}Step 8: Standardized Tool Execution via Model Context Protocol}
\vspace{0.15em}

\begin{algolist}
  \item[8.1] Interface agents to external tools via Model Context Protocol (MCP):
    \begin{algosublist}
      \item[8.1.1] Standardize client-server communication using JSON-RPC protocol specifications.
      \item[8.1.2] Enable dynamic tool discovery: inspect exposed tool schemas, input parameters, and return types.
      \item[8.1.3] Decouple agent reasoning logic from external tool implementation details.
    \end{algosublist}
  \item[8.2] Orchestrate multi-domain tool execution ecosystem:
    \begin{algosublist}
      \item[8.2.1] \textbf{System \& Dev Tools:} Local Filesystem (file I/O), Sandboxed CLI (isolated shell), GitHub API.
      \item[8.2.2] \textbf{Productivity Tools:} Google Workspace MCP (Gmail messaging, Google Drive document access).
      \item[8.2.3] \textbf{Web Intelligence:} Browser MCP for real-time document search, retrieval, and scraping.
    \end{algosublist}
  \item[8.3] Enforce Human-in-the-Loop (HITL) authorization policy:
    \begin{algosublist}
      \item[8.3.1] Permit autonomous execution for non-destructive, read-only tool queries.
      \item[8.3.2] Intercept state-mutating actions (file modification, email dispatch, CLI execution) for user approval.
    \end{algosublist}
\end{algolist}
\end{frame}

% -----------------------------------------------------------------------------
% STEP 9: CONTEXT COMPACTION & SHARED MEMORY
% -----------------------------------------------------------------------------
\begin{frame}[t]{Algorithm: Step 9 -- Context Compaction and Shared Memory}
\vspace{-0.15em}
{\large\bfseries\color{lavenderDark}Step 9: Context Compaction and Shared Memory Management}
\vspace{0.15em}

\begin{algolist}
  \item[9.1] Address context window constraints and token consumption limits:
    \begin{algosublist}
      \item[9.1.1] Mitigate fixed LLM context window boundaries to eliminate token overflow and latency degradation.
      \item[9.1.2] Prevent quadratic attention computation overhead across long multi-agent interaction sessions.
    \end{algosublist}
  \item[9.2] Execute dynamic context compaction and artifact pruning:
    \vspace{-0.2em}
    $$C_{t+1} = \mathrm{Prune}\big(C_t \cup \Delta \text{Artifacts} \setminus \text{Stale Tokens}\big)$$
    \vspace{-0.2em}
    {\footnotesize\color{prismDark}%
      \textbf{\color{lavenderDark}Where:}\enspace $C_{t+1}, C_t$: Next and current context buffers $\;\vert\;$ $\Delta\text{Artifacts}$: Verified upstream deliverables $\mathcal{O}_i$ $\;\vert\;$ $\text{Stale Tokens}$: Discarded scratchpad tokens \& raw tool logs $\;\vert\;$ $\mathrm{Prune}(\cdot)$: Heuristic eviction filter.}
    \vspace{0.15em}
    \begin{algosublist}
      \item[9.2.1] Forward only verified parent output artifacts $\mathcal{O}_i$ to dependent downstream child nodes.
      \item[9.2.2] Evict intermediate reasoning scratchpads, redundant tool telemetry, and discarded draft tokens.
    \end{algosublist}
  \item[9.3] Persist global execution state using client-side database (\texttt{PrismDB}):
    \begin{algosublist}
      \item[9.3.1] Commit session turns, task graphs, and tool execution logs to browser IndexedDB via Dexie.js.
      \item[9.3.2] Ensure client data privacy, full offline execution persistence, and instant crash recovery.
    \end{algosublist}
\end{algolist}
\end{frame}

% -----------------------------------------------------------------------------
% STEP 10: ORPO PREFERENCE MODEL & QUALITY ASSERTION
% -----------------------------------------------------------------------------
\begin{frame}[t]{Algorithm: Step 10 -- ORPO Preference Model \& Verification}
\vspace{-0.15em}
{\large\bfseries\color{lavenderDark}Step 10: Single-Stage Preference Alignment (ORPO) and Verification}
\vspace{0.15em}

\begin{algolist}
  \item[10.1] Single-stage reference-free preference alignment formulation (ORPO):
    \vspace{-0.25em}
    $$\mathcal{L}_{\mathrm{ORPO}} = \mathcal{L}_{\mathrm{SFT}} + \lambda_{\mathrm{OR}} \cdot \mathcal{L}_{\mathrm{OR}} \quad \text{where} \quad \mathcal{L}_{\mathrm{OR}} = -\log \sigma \left( \log \left[ \frac{\mathrm{odds}_\theta(y_w \mid x)}{\mathrm{odds}_\theta(y_l \mid x)} \right] \right)$$
    \vspace{-0.2em}
    {\footnotesize\color{prismDark}%
      \textbf{\color{lavenderDark}Where:}\enspace $\mathcal{L}_{\mathrm{ORPO}}$: Total loss $\;\vert\;$ $\mathcal{L}_{\mathrm{SFT}}$: SFT cross-entropy loss $\;\vert\;$ $\lambda_{\mathrm{OR}}$: Odds weight $\;\vert\;$ $\mathcal{L}_{\mathrm{OR}}$: Odds ratio penalty $\;\vert\;$ $\mathrm{odds}_\theta(y \mid x) = \frac{P_\theta(y \mid x)}{1 - P_\theta(y \mid x)}$: Output sequence odds $\;\vert\;$ $y_w, y_l$: Chosen / rejected responses for prompt $x$.}
    \vspace{0.15em}
    \begin{algosublist}
      \item[10.1.1] Jointly maximize chosen likelihood while penalizing rejected tokens ($2\times$ memory savings).
      \item[10.1.2] Supervise via UltraFeedback pairs $\mathcal{D} = \{(x, y_w, y_l)\}$, bypassing unstable RLHF loops.
    \end{algosublist}
  \item[10.2] Enforce runtime quality gating via preference score barrier $\mathcal{S}_{\mathrm{ORPO}}(x, y) \ge \tau_{\mathrm{approval}}$:
    \begin{algosublist}
      \item[10.2.1] Audit deliverables using fine-tuned model (\texttt{reward\_orpo\_v2}, Qwen2.5-1.5B LoRA).
      \item[10.2.2] Score candidate outputs on instruction adherence, structural constraints, and safety guardrails.
    \end{algosublist}
  \item[10.3] Handle quality failures via targeted DAG re-planning:
    \begin{algosublist}
      \item[10.3.1] Capture diagnostic error traces and schema deviations when output falls below $\tau_{\mathrm{approval}}$.
      \item[10.3.2] Dispatch error logs to DAG Planner for targeted branch re-planning without system abort.
    \end{algosublist}
\end{algolist}
\end{frame}

% -----------------------------------------------------------------------------
% STEP 11: FINAL DELIVERABLE SYNTHESIS & COMMIT
% -----------------------------------------------------------------------------
\begin{frame}[t]{Algorithm: Step 11 -- Final Deliverable Synthesis and Commit}
\vspace{-0.15em}
{\large\bfseries\color{lavenderDark}Step 11: Artifact Assembly and Global State Commit}
\vspace{0.15em}

\begin{algolist}
  \item[11.1] Aggregate and assemble multi-agent deliverables:
    \begin{algosublist}
      \item[11.1.1] Harvest verified partial output artifacts $\mathcal{O}_i$ from all terminal nodes in the completed DAG.
      \item[11.1.2] Coalesce research summaries, verified code diffs, and formatted text into unified final deliverable.
    \end{algosublist}
  \item[11.2] Calculate end-to-end composite delivery quality metric:
    \vspace{-0.2em}
    $$\text{Delivery Quality } D = \text{Correctness} \times \text{Completeness} \times \text{Timeliness}$$
    \vspace{-0.2em}
    {\footnotesize\color{prismDark}%
      \textbf{\color{lavenderDark}Where:}\enspace $D \in [0, 1]$: Composite quality metric $\;\vert\;$ \text{Correctness}: Schema conformance \& execution pass rate $\;\vert\;$ \text{Completeness}: Fulfillment of goal constraints $\;\vert\;$ \text{Timeliness}: Latency penalty score.}
    \vspace{0.15em}
    \begin{algosublist}
      \item[11.2.1] Apply multiplicative scoring penalty: failure in any single dimension sharply penalizes overall quality.
      \item[11.2.2] Log quality vector to persistent performance ledger to calibrate future provider routing models.
    \end{algosublist}
  \item[11.3] Commit finalized session state with ACID guarantees:
    \begin{algosublist}
      \item[11.3.1] Atomically write completed conversation records, DAG execution graphs, and tool receipts to \texttt{PrismDB}.
      \item[11.3.2] Record end-to-end telemetry: aggregate latency, per-agent token consumption, and audit receipts.
    \end{algosublist}
\end{algolist}
\end{frame}

% -----------------------------------------------------------------------------
% STEP 12: OUTPUT DELIVERY & REACTIVE WORKSPACE UI
% -----------------------------------------------------------------------------
\begin{frame}[t]{Algorithm: Step 12 -- Output Delivery and User Display}
\vspace{-0.15em}
{\large\bfseries\color{lavenderDark}Step 12: Output Delivery and Reactive Workspace Presentation}
\vspace{0.15em}

\begin{algolist}
  \item[12.1] Synthesize and package final deliverable artifacts:
    \begin{algosublist}
      \item[12.1.1] Compile structured digital artifacts (Markdown reports, verified code patches, email/document drafts).
      \item[12.1.2] Generate comprehensive audit receipt containing agent provenance, execution duration, and tool logs.
    \end{algosublist}
  \item[12.2] Render presentation layer via multi-pane workspace interface:
    \begin{algosublist}
      \item[12.2.1] Update interactive DAG HUD with real-time node states (\texttt{Queued} $\to$ \texttt{Running} $\to$ \texttt{Success} / \texttt{Failed}).
      \item[12.2.2] Mount Artifact Viewer canvas providing rich Markdown rendering, syntax highlighting, and export controls.
      \item[12.2.3] Present interactive approval cards for pending external actions with explicit authorization controls.
    \end{algosublist}
  \item[12.3] Synchronize client state reactively via local database transactions:
    \begin{algosublist}
      \item[12.3.1] Commit verified state transitions and output artifacts directly into browser IndexedDB (\texttt{PrismDB}).
      \item[12.3.2] Trigger Dexie \texttt{useLiveQuery()} observers to update UI components reactively upon local database commits.
    \end{algosublist}
\end{algolist}
\end{frame}

\end{document}
